\subsection{Inference Pipeline}
\label{sec:inference}

The trained PMDF-Net operates as a feed-forward denoising module requiring no iterative optimization, averaging, or per-specimen fine-tuning at inference time. Given a raw or low-averaged input signal of length $L$ sampled at the same rate as training data, the network produces a denoised waveform of identical length through a single forward pass. This contrasts directly with conventional signal averaging, which demands $N$ repeated acquisitions summed point-wise to achieve an SNR improvement of $10\log_{10}(N)$ dB.

\textbf{Input and output format.} The inference input is a single ultrasonic waveform---either a single-pulse acquisition or an $N$-shot averaged signal where $N \in \{1, 2, 4, 8\}$, matching the training input regime. No additional metadata, reference signals, or calibration acquisitions are required. The output is a point-wise denoised waveform $\hat{y} \in \mathbb{R}^L$, ready for subsequent arrival-time picking, amplitude extraction, or defect detection pipelines.

\textbf{Computational cost.} Inference throughput depends on hardware platform and signal length. On a desktop GPU (NVIDIA RTX 3080), a forward pass for a 2048-point signal completes in approximately 2--5 ms, yielding a throughput of 200--500 signals per second in batch mode. On a mid-range edge GPU (NVIDIA Jetson Orin), latency increases to 15--30 ms per signal, sufficient for real-time processing at 30--60 signals per second. A CPU-only inference path is also supported for deployment on standard workstations, achieving 50--100 signals per second for the same signal length.

\textbf{Speedup relative to averaging.} Consider the practical scenario of achieving an effective SNR equivalent to 64-shot averaging. A single forward pass of PMDF-Net on a 1-shot input achieves comparable denoising quality in approximately 3--5 ms on GPU, whereas 64-shot hardware averaging requires 64 times the acquisition duration plus summation overhead. Even accounting for network inference, the proposed approach delivers a $10\times$ to $100\times$ speedup in effective acquisition-to-output latency, depending on the $N$-shot baseline under comparison.

\textbf{Real-time deployment.} The fully convolutional architecture of PMDF-Net contains no recurrent connections or temporal dependencies that require sequential processing, enabling trivial parallelization across signals. Model quantization (INT8) reduces the footprint to under 50 MB, making GPU-edge deployment feasible for in-situ inspection systems. No fine-tuning or domain adaptation is required when applying the network to new specimens: the model generalizes from intact-aluminum training data to defect-containing aluminum and CFRP targets without parameter updates, as validated in Section~\ref{sec:experiments}.
